\section{Formulation}
\label{sec:problem}

A branching step is the object \sys admits.
\S\ref{sec:bg-full} is the reason prefix caching, disaggregation, and decode-time pruning each pay for a non-winner before that admission can run;
\Cref{tab:trad,tab:notation} record the stage and the symbols.

\begin{definition}[Branching step]
\label{def:step}
A parent $u$ holds a committed trunk $x_u$ with $|x_u|{=}L$.
A fan-out is a set of residuals $\{\rho_i\}_{i=0}^{k-1}$, $|\rho_i|{=}\ell_i$.
The bound winner $\star$ leaves the spine $x_\star=x_u\Vert\rho_\star$, $|x_\star|{=}L{+}\ell_\star$.
\end{definition}

Let $b$ be KV bytes per token.
The three occupancies of one step are
\begin{align}
M_{\mathrm{clone}} &= b\bigl(kL+\textstyle\sum_i\ell_i\bigr), \label{eq:mem}\\
M_{\mathrm{CoW}} &= b\bigl(L+\textstyle\sum_i\ell_i\bigr), \nonumber\\
M_{\mathrm{spine}} &= b\bigl(L+\ell_\star\bigr). \nonumber
\end{align}
$M_{\mathrm{clone}}$ is the first-expand footprint of recompute, APC, and baseline disaggregation.
$M_{\mathrm{CoW}}$ is one shared trunk with every residual live.
$M_{\mathrm{spine}}$ is the footprint after abort.
On a cold radix the same span is what disaggregation transfers,
\begin{equation}
\label{eq:xfer-stock}
\Xfer_{\mathrm{PD}}=\Work_{\mathrm{APC}}=kL+\textstyle\sum_i\ell_i,
\end{equation}
because a block hash exists only after the tokens do, and the connector moves every page the prefiller computed~\cite{vllm,distserve,splitwise}.
$\Work$ is the token count issued to a prefill kernel before $\star$ is known.
$\Xfer$ is the token count inserted on that connector.
The admission rule below replaces both with the miss tail of the residuals it keeps, so $L$ is absent from $\Work$ and from $\Xfer$.

\subsection{Where the Fan-out Is Paid}
\label{sec:traditional}

Four mechanisms already sit on this step, and each one can reject a child only after a fixed stage (\Cref{tab:trad}).
Recompute clones the trunk.
APC and a radix index share a block only after that block has been published, so the first fan-out in a batch still materializes $k$ copies of $L$; a later replay of the winner is a hit.
Prefill--decode disaggregation transfers the pages the prefiller has already computed, aborted residuals included, which is why $\Xfer_{\mathrm{PD}}=\Work_{\mathrm{APC}}$ on a cold index.
ESC, Speculative Rejection, and DPTS score a generated prefix, so the residual prefill, and under disaggregation the transfer, precede the drop.
\sys aliases one trunk, withholds a repeated loop before the kernel, and continues every other low score only through the probe in \eqref{eq:dec}.
The connector then carries the miss tail of $\Keep$, not the $k$-way clone.

\begin{table*}[t]
\centering
\caption{Where each method pays for one fan-out of $k$ residuals on a trunk of length $L$.
``First expand'' is the GPU prefill issued before any child has published a block.
``Rejection'' is the earliest point at which a non-winner can leave.
The \sys row is the action of \S\ref{sec:app}.}
\label{tab:trad}
\scriptsize
\setlength{\tabcolsep}{3.5pt}
\begin{tabular}{@{}p{0.18\textwidth}p{0.22\textwidth}p{0.30\textwidth}p{0.22\textwidth}@{}}
\toprule
Method & First expand & Rejection & Connector \\
\midrule
Recompute & $k$ clones of $L$ & end of the request & none \\
APC, radix, hash prefill & $k$ clones of $L$ & LRU after a block is published & none; a later replay of the winner is a hit \\
P/D disaggregation & $k$ clones of $L$ & after \texttt{insert} & $kL+\sum_i\ell_i$ \\
ESC, Speculative Rejection, DPTS & every residual & after a generated prefix & KV of every prefilled child \\
\sys & one alias of $L$ & loop before the kernel; other low scores after a $t_0$ probe & miss tail of $\Keep$ \\
\bottomrule
\end{tabular}
\end{table*}
